\section{收益与现金分配的正确口径}

本报告不为个股、ETF、利率债、黄金或组合给出点状年化收益、固定票息或固定现金分配预测。资产管理产品不能承诺本金和回报；ETF分配也来自基金资产而非额外财富。以某一期历史分配率、某日价格或“股息率加增长率”直接合成组合收益，无法反映价格、估值、税费、费用、再投资时点和分配变化。

\section{长期规划敏感性，不是预测}

下表只回答“若长期名义总回报参数$r$落在不同规划输入，十年公式结果会怎样”，而不判断哪个参数最可能发生。$r$是扣除税费、交易成本和基金费用之前的假设总回报输入；不含年度追加、取现、再投资时点或路径，也不生成实际现金$D_0$。

\begin{exhibitbox}[五年以上规划参数]
\small
\begin{tabularx}{\exhibitboxwidth}{L{3cm}P P}
\toprule
\textbf{年化总回报参数} & \textbf{在本报告中的用途} & \textbf{明确不代表} \\
\midrule
3.0\% & 保守端的长期复利公式敏感性输入 & 下一年度收益、利息或股息承诺。 \\
4.5\% & 中间的长期复利公式敏感性输入 & “基准预测”或可提取现金。 \\
6.0\% & 较高端的长期复利公式敏感性输入 & 目标收益、估值上行或产品分配保证。 \\
\bottomrule
\end{tabularx}
\end{exhibitbox}

规划区间为年化3\%--6\%，适用期限为五年以上。它只用于比较期限、损失预算和复利算术，不替代执行日前的产品和市场复核。

\section{实际现金分配的记录方法}

季度复核时，定义

\[
D_0=\sum \text{过去12个月实际到账的利息、分配和股息（税后、扣费后）}.
\]

$D_0$ 是已经发生的现金记录，不外推为未来分配率。每一笔记录须包含到账日、来源资产层、税前金额、税额、费用、净额和结算或证据凭证。若实际分配被削减，使用 $D_h=D_0(1-h)$，其中 $h\in\{0,25\%,50\%,100\%\}$，仅评估可再投资现金减少对复利速度的影响；只有已有实际$D_0$记录时才计算该敏感性。任何分配为零的年份也不改变组合风险上限；资金可以留在组合内受保障流动性，不能以杠杆或追高补回。

\section{不纳入本模型的数字}

\begin{itemize}
  \item 未来一年的组合回报、股息率、债券ETF票息、货币工具收益或黄金收益；
  \item 每年可提取的固定现金、生活费覆盖额、税后收入或“收益来源占比”；
  \item 基于观察池的个股等权收益、资本利得假设或目标价格；
  \item 未经执行日核验的ETF价格、净值、折溢价、成交量、价差和交易成本。
\end{itemize}

\qualitynote{长期公式敏感性与实际现金台账彼此独立：$r$不预测$D_0$，$D_0$也不改变已披露的公式案例。}
